\subsubsection{BSuite tasks}
\label{sec:app_bsuite_tasks}

We provide a brief description below of the three behavior suite (BSuite) reinforcement learning tasks in \citet{osband2019behaviour}:
\begin{enumerate}
    \item Catch: A 10x5 Tetris-grid with single block falling per column. The agent can move left/right in the bottom row to ‘catch’ the block. Illustrated in Figure~\ref{fig:bsuite_catch}.
    \item Mountain Car: The agent drives an underpowered car up a hill~\citep{moore1990efficient}. Illustrated in Figure~\ref{fig:bsuite_mountain_car}.
    \item Cartpole: The agent can move a cart left/right on a plane to keep a balanced pole upright~\citep{barto1983neuronlike}. Illustrated in Figure~\ref{fig:bsuite_cartpole}.
\end{enumerate}

\begin{figure}[tbph!]
     \centering
     \begin{subfigure}[b]{0.3\textwidth}
         \centering
         \includegraphics[width=\textwidth]{Appendix/Catch.png}
         \caption{Catch}
         \label{fig:bsuite_catch}
     \end{subfigure}
     \hfill
     \begin{subfigure}[b]{0.3\textwidth}
         \centering
         \includegraphics[width=\textwidth]{Appendix/Mountain_Car.png}
         \caption{Mountain Car}
         \label{fig:bsuite_mountain_car}
     \end{subfigure}
     \hfill
     \begin{subfigure}[b]{0.3\textwidth}
         \centering
         \includegraphics[width=\textwidth]{Appendix/Cartpole.png}
         \caption{Cartpole}
         \label{fig:bsuite_cartpole}
     \end{subfigure}
        \caption{Three BSuite tasks. Figures are from \citet{osband2019behaviour}.}
        \label{fig:bsuite}
\end{figure}

All tasks have a real-valued state space $\mathcal{S}\subseteq \mathcal{R}^{D}$ and a discrete action space $\mathcal{A}=\{0, 1, 2, \dots, A-1\}$. The state dimension and number of actions are provided in Table~\ref{tab:bsuite_dims}.

\begin{table}
    \centering
    \begin{tabular}{lccc}
         &  \textbf{Catch} &\textbf{Mountain Car} & \textbf{Cartpole} \\ \hline
        $D$ & 50 & 3 & 6\\
        $A$ & 3 & 3 & 3
    \end{tabular}
    \caption{Dimensions of BSuite tasks}
    \label{tab:bsuite_dims}
\end{table}

\subsubsection{Fitted Q Evaluation}
\label{sec:app_fqe}
Fitted Q Evaluation (FQE)~\citep{le2019batch} is a simple variant of the Fitted Q Iteration (FQI) algorithm~\citep{ernst2005tree}. Instead of learning the Q function of an optimal policy as FQI, FQE estimates the Q function of a fixed policy $\pi$. It is an iterative algorithm with randomly initialized the Q function parameters $\theta_0$. At iteration $k\geq 1$ with the current estimate of Q function $Q^{\pi}(s, a|\theta_{k-1})$, it updates parameters $\theta_k$ by solving the following regression problem using a least squares approach:

\begin{align}\label{eq:fqe}
    r + \gamma Q^{\pi}(s',a'|\theta_{k-1}) =& Q^{\pi}(s,a|\theta_{k}) \\
       &+ \gamma \left(Q^{\pi}(s',a'|\theta_{k-1})- \expect[s'\sim P(\cdot|s,a), a'\sim \pi(\cdot | s')]{Q^{\pi}(s',a'|\theta_{k-1})}\right) \nonumber \\
       &+r-\expect[r\sim R(\cdot | s, a,s')]{r} \nonumber,
\end{align}
where the regression function to estimate is $Q^{\pi}(s,a|\theta_{k})$, the observed outcome is the term on the LHS of \eqref{eq:fqe} and the residual is the sum of terms in the second and third lines of  \eqref{eq:fqe}. Comparing the equation above with (\ref{eq:SCMforRL}), FQE reformulates the regression problem and moves the confounded part in the treatment of (\ref{eq:SCMforRL}), that is $\gamma Q^\pi(s', a')$, to the outcome. The regression problem at each iteration is therefore not confounded any more. If FQE converges, it finds a solution of the following equation
\begin{align}
    Q^{\pi}(s,a|\theta) = P(\expect[r|s,a]{r} + \gamma \expect[s', a'|s,a]{Q^{\pi}(s',a'|\theta)})\,,
\end{align}
where $P(\cdot)$ is the $L_2$ projection operator that maps a function of $(s, a)$ onto the (parameterized) $Q$ function space.